% !TEX program = lualatex
\documentclass[aspectratio=43,professionalfonts,handout]{beamer}
\usepackage{beamer_template}

\newcommand{\LectureNum}{06}
\newcommand{\LectureTitle}{２次関数と２次方程式（２）}
\newcommand{\TextPages}{pp.\,83, 86-87}
\newcommand{\NextTopic}{２次関数と２次不等式}
\newcommand{\NextPages}{pp.\,84-85}
\newcommand{\CourseName}{基礎数学B}
\renewcommand{\TermName}{前期}

\begin{document}

% -----------------------------------------------
% スライド1: 表紙＋目標＋用語・定義
% -----------------------------------------------
\begin{frame}{\CourseName\quad 第\LectureNum 回「\LectureTitle 」}
  {\small \TermName\quad \TeacherName\quad ／\quad 教科書 \TextPages}

  \begin{block}{用語・定義}
  \begin{description}[labelwidth=5.5em]
    \item[\kw{前回の復習}] $D=b^2-4ac$ の符号で放物線と $x$ 軸の位置関係が決まる
    \item[\kw{常に $y>0$}] $a>0$ かつ $D<0$（放物線が $x$ 軸より常に上）
    \item[\kw{常に $y\geq 0$}] $a>0$ かつ $D\leq 0$（$x$ 軸と接するか常に上）
  \end{description}
  \end{block}
\end{frame}

% -----------------------------------------------
% スライド2: 公式・性質
% -----------------------------------------------
\begin{frame}{公式・性質}
  \begin{block}{}
  \begin{itemize}
    \item $a>0$ かつ常に $y>0$ $\Longleftrightarrow$ $D<0$
    \item $a>0$ かつ常に $y\geq 0$ $\Longleftrightarrow$ $D\leq 0$
    \item 解と係数の関係（参考）：$\alpha+\beta=-\tfrac{b}{a}$、$\alpha\beta=\tfrac{c}{a}$
  \end{itemize}
  \end{block}
\end{frame}

% -----------------------------------------------
% スライド3: 例1→問1
% -----------------------------------------------
\begin{frame}{例1→問1}
  \begin{exampleblock}{練習1-A 問6) p.86（平行移動して原点を通る）}
    （板書で解説）
  \end{exampleblock}

  \vfill

  \begin{block}{練習1-A 問7) p.86\quad 自分で解いてみよう}
    （ノートに解くこと）
  \end{block}
\end{frame}

% -----------------------------------------------
% スライド4: 例2→問2
% -----------------------------------------------
\begin{frame}{例2→問2}
  \begin{exampleblock}{練習1-B 問4) p.87（内接する長方形）}
    （板書で解説）
  \end{exampleblock}

  \vfill

  \begin{block}{練習1-B 問5) p.87\quad 自分で解いてみよう}
    （ノートに解くこと）
  \end{block}
\end{frame}

% -----------------------------------------------
% まとめ
% -----------------------------------------------
\begin{frame}{まとめ}
  \begin{itemize}
    \item 判別式と $x$ 軸の位置関係の3パターンを確実に使い分ける
    \item 放物線が $x$ 軸より上・下にある条件も判別式で判定
  \end{itemize}

  \vfill
  {\small 基本問題：154, 163}
\end{frame}

\end{document}
